% The obligations of process steps. From map block p-obls, deadlock-free
% reading. The disjointness side conditions of Q-Tau, Q-Send and Q-Recv
% (C' disjoint from chans Sigma, ech[dom Theta_2] disjoint from C, the
% first conjuncts of send_split and recv_ok) are covered by the convention
% on disjoint contexts. send_split, recv_ok, recv_back and recv_env are
% opened, matches, closed and iface of def:proc-conds.
\begin{figure}[tbp]
  \small
  \judghead{Obligations of process steps}{\mathsf{Pobl}^{W}_{ev}(\proc')}
  \begin{rules}
    \infer*[lab=Q-Tau]
      {\proc'\ \mathsf{at}\ \senv}
      {\mathsf{Pobl}^{W}_{\tau}(\proc')}
    \label{rule:Q-Tau}
    \and
    \infer*[lab=Q-Send]
      {\senv(o) \doteq (\sSend{\bar\iota}{\tyB}.\sty)_{x} \\
       \senv \setminus o = \senv_v \uplus \senv_r \\\\
       \mathrm{dom}\,\cTheta_e = \{o' \in \fends(\valB) \mid \ech(o') \in C\} \\
       \mathsf{opened}(\senv_e, \cTheta_e) \\\\
       \valB \in \Vrel{W}{\ix}{\senv_v \cup \cTheta_e}{\payload{\tyB}{x}} \\
       \proc'\ \mathsf{at}\ (\senv_r \cup \senv_e)[o \mapsto \sty_{\pplus{x}}]}
      {\mathsf{Pobl}^{W}_{\tsend\,o\,\valB}(\proc')}
    \label{rule:Q-Send}
    \and
    \infer*[lab=Q-Recv]
      {\senv(o) \doteq (\sRecvT{\bar\iota}{\tyB}.\sty)_{x} \\
       \senv' = \senv[o \mapsto \sty_{\pplus{x}}] \\\\
       \forall \ixs < \ix.\ \forall \valB, \cTheta_2. \\\\
       \mathsf{matches}(\senv', \cTheta_2) \land \valB \in \Vrel{W}{\ixs}{\cTheta_2}{\payload{\tyB}{x}} \\\\
       \Rightarrow\ (C' \cup \mathsf{closed}(\senv', \cTheta_2),\ E') \in \Prel{W}{\ixs}{\mathsf{iface}(\senv', \cTheta_2)}{\ty}}
      {\mathsf{Pobl}^{W}_{\trecv\,o\,\valB}(\proc')}
    \label{rule:Q-Recv}
  \end{rules}
  \figrule
  \Description{Inference rules for the obligations of an internal step, of a send to the environment and of a receive from the environment.}
  \caption{Obligations of process steps for $\proc = (C, E)$ and $\proc' = (C', E')$, under \cref{def:rel-anatomy,def:proc-conds}. $\tau$ is an internal step, a message between two threads of $\proc$ included, and a close, a choice and a link on an interface end repeat O-Close, O-Select, O-Branch and O-Link with $\Sigma$ in place of $\Theta$ (\cref{app:p-obls}).}
  \label{fig:proc-obl}
\end{figure}
